% --- Perfil ---
\section{Perfil}
\begin{onecolentry}
	Soy graduado en Ingeniería de Telecomunicación, habiendo cursado el grado y el máster en la Universidad de Alcalá. Dentro del campo de la ingeniería, estoy especialmente interesado en la investigación aplicada, los sistemas distribuidos y el procesamiento intensivo, habiendo orientado mi trayectoria profesional hacia el desarrollo de plataformas de análisis de datos, la computación multihilo y en clúster, y la construcción de entornos reproducibles para experimentación. Mi experiencia combina trabajo científico y desarrollo técnico, incluyendo automatización de despliegues, procesamiento paralelo y uso de tecnologías de virtualización, contenedores y orquestación para acelerar ciclos de prueba, facilitar la validación de resultados y dar soporte a infraestructuras de investigación.
\end{onecolentry}

\vspace{-5pt}

% --- Educación ---
\section{Educación}

\begin{onecolentry}
	\begin{tabularx}{\textwidth}{@{}>{\raggedright\arraybackslash}X>{\raggedright\arraybackslash}p{3.6cm}>{\raggedleft\arraybackslash}p{2.4cm}@{}}
		\textbf{Máster en Ingeniería en Tecnologías de Telecomunicación} & Universidad de Alcalá & \textit{Febrero 2026}
	\end{tabularx}
	
	\vspace{0.05 cm}
	\hspace{0.2cm}Premio extraordinario
\end{onecolentry}

\vspace{0.2 cm}

\begin{onecolentry}
	\begin{tabularx}{\textwidth}{@{}>{\raggedright\arraybackslash}X>{\raggedright\arraybackslash}p{3.6cm}>{\raggedleft\arraybackslash}p{2.4cm}@{}}
		\textbf{Grado en Ingeniería en Tecnologías de Telecomunicación} & Universidad de Alcalá & \textit{Julio 2024}
	\end{tabularx}
	
	\vspace{0.05 cm}
	\hspace{0.2cm}2.º puesto de promoción y premio extraordinario
\end{onecolentry}

\vspace{-7pt}

% --- Experiencia ---
\section{Experiencia}

\begin{twocolentry}{\textit{Septiembre 2024 – Enero 2026}}
	\textbf{Investigación en servicios e infraestructuras de red} \\
	Universidad de Alcalá, Alcalá de Henares
\end{twocolentry}
\begin{onecolentry}
	\begin{highlights}
		\item Diseñé y desplegué servicios en contenedores, utilizando Docker y containerd, así como orquestadores como Docker Compose, Kubernetes y OpenShift.
		\item Integré flujos de despliegue automatizados en entornos empresariales mediante GitHub Workflows y Vagrant, tanto para entornos de prueba como de producción.
		\item Desarrollé pipelines de procesamiento paralelo en Python, MATLAB y C para cargas intensivas y análisis de datos.
	\end{highlights}
\end{onecolentry}

\vspace{0.2 cm}

\begin{twocolentry}{\textit{Febrero 2024 – Agosto 2024}}
	\textbf{Sistemas, seguridad y soporte técnico} \\
	Correos y Telégrafos S.A., Madrid
\end{twocolentry}
\begin{onecolentry}
	\begin{highlights}
		\item Trabajé en la operación y documentación de equipos de red y sistemas corporativos.
		\item Gestioné servicios de seguridad y autenticación como RADIUS, LDAP, OAuth2 y TLS.
		\item Participé en despliegues y prácticas DevOps, integrando automatización en pipelines de CI/CD.
		\item Administré sistemas Linux empresariales, principalmente RHEL y Amazon Linux 2.
	\end{highlights}
\end{onecolentry}

\vspace{0.2 cm}

\begin{twocolentry}{\textit{Enero 2023 – Enero 2024}}
	\textbf{Ingeniería e investigación aplicada} \\
	Universidad de Alcalá / Correos y Telégrafos S.A., Alcalá de Henares
\end{twocolentry}
\begin{onecolentry}
	\begin{highlights}
		\item Desarrollé plataformas web de análisis de datos, basadas en arquitectura MVC, para proyectos de investigación.
		\item Implementé soluciones de computación multihilo y en clúster para tareas de procesamiento intensivo.
		\item Me encargué de la administración de entornos Linux y de sistemas de computación en red.
		\item Colaboré en la actividad científica del grupo y participé como coautor en publicaciones \href{https://orcid.org/0009-0000-6862-1872}{ORCID}.
	\end{highlights}
\end{onecolentry}

\vspace{-7pt}

% --- Aptitudes ---
\section{Aptitudes}

\begin{onecolentry}
	\textbf{Programación:} C/C++, Python, MATLAB, Bash, Groovy
\end{onecolentry}
\begin{onecolentry}
	\textbf{Procesamiento y computación:} Computación multihilo, procesamiento paralelo, computación en clúster
\end{onecolentry}
\begin{onecolentry}
	\textbf{Infraestructura:} Linux, Docker, containerd, Docker Compose, Kubernetes, Red Hat OCP, Vagrant
\end{onecolentry}
\begin{onecolentry}
	\textbf{CI/CD:} GitHub Workflows, Jenkins, GitHub/GitLab
\end{onecolentry}